\einheitenkopf{Einheit 3 von 5 · Streifen- und Kreisdiagramm}
\erg{22}{a) $^\circ$ \quad b) \% \quad c) $^\circ$ \quad d) \% \quad e) $^\circ$}
\erg{23}{a) $4\,$cm \quad b) $2\,$cm \quad c) $7\,$cm \quad d) $0{,}5\,$cm \quad e) $4{,}5\,$cm \quad f) Rad $2{,}5\,$cm, Bus $4\,$cm, zu Fuß $2\,$cm, Auto $1{,}5\,$cm (zusammen $10\,$cm) \quad g) $100-75=25\,\%$}
\erg{24}{a) $180^\circ$ \quad b) $90^\circ$ \quad c) $36^\circ$ \quad d) $270^\circ$ \quad e) $0{,}12 \cdot 360^\circ = 43{,}2^\circ$ \quad f) $135^\circ$ \quad g) $8:25 = 0{,}32 \rightarrow 115{,}2^\circ$ \quad h) $15:400 = 0{,}0375 \rightarrow 13{,}5^\circ$ \quad i) $3:7 \approx 0{,}4286 \rightarrow 154{,}3^\circ$}
\erg{25}{a) und b) Sektoren von $90^\circ$ und $120^\circ$, vom eingezeichneten Radius aus im Uhrzeigersinn \quad c) $360^\circ - 90^\circ - 120^\circ = 150^\circ$ \quad d) Fußball $\tfrac{100}{240}\cdot 360^\circ = 150^\circ$, Handball $120^\circ$, Judo $90^\circ$ \quad e) Kreis mit drei Sektoren $150^\circ$, $120^\circ$, $90^\circ$, beschriftet}
\erg{26}{a) größter Sektor Gas $38\,\%$, dann Öl $27\,\%$, dann Fernwärme $21\,\%$, kleinster Wärmepumpe $14\,\%$ \quad b) Öl mit $27\,\%$ \quad c) $17 : 450 \approx 0{,}0378 \rightarrow 0{,}0378 \cdot 360^\circ \approx 13{,}6^\circ$}
\erg{27}{a) Ben hat die Prozentzahl als Gradzahl eingetragen. Der Anteil muss erst mit $360^\circ$ multipliziert werden: $0{,}24 \cdot 360^\circ = 86{,}4^\circ$. \quad b) $0{,}45 \cdot 360^\circ = 162^\circ$}
\erg{28}{a) Der ganze Kreis sind $360^\circ$ und $100\,\%$. Ein Prozent ist ein Hundertstel davon: $360^\circ : 100 = 3{,}6^\circ$. \quad b) Sein Anteil ist größer als $50\,\%$, denn der halbe Kreis entspricht genau der Hälfte des Ganzen.}
